\documentclass[12pt,a4paper]{article}

\usepackage{fontspec}
\usepackage[spanish]{babel}
\usepackage{geometry}
\usepackage{setspace}
\usepackage{ragged2e}

\geometry{
	left=3cm,
	right=2.5cm,
	top=2.5cm,
	bottom=2.5cm
}

\setmainfont{Times New Roman}

\begin{document}
	
	\onehalfspacing
	\justifying
	\setlength{\parindent}{0pt}
	\setlength{\parskip}{0.8em}
	
	\begin{center}
		{\Large \textbf{Dependencia de recursos naturales y trayectorias de transición estructural (1980–2022)}}\\
		\vspace{0.4cm}
		Julián Alberto Delgadillo Marín\\
		Maestría en Economía Aplicada – UBA\\
		Proyecto de Trabajo Final de Maestría
	\end{center}
	
	\vspace{0.8cm}
	
	\vspace{0.8cm}
	
	La abundancia de recursos naturales no renovables ha sido históricamente considerada una ventaja comparativa potencial para el crecimiento económico; sin embargo, la evidencia comparada muestra que las economías intensivas en recursos extractivos exhiben trayectorias divergentes en términos de diversificación productiva, acumulación de capacidades e innovación. Mientras algunos países han logrado transformar las rentas derivadas de la explotación de hidrocarburos y minerales en procesos sostenidos de transición estructural, otros permanecen atrapados en configuraciones productivas concentradas, vulnerables a la volatilidad externa y con limitada complejidad económica. Esta heterogeneidad constituye un problema central para la economía del desarrollo contemporánea, en la medida en que cuestiona tanto las visiones deterministas de la denominada “maldición de los recursos” como las perspectivas excesivamente optimistas que asumen que la abundancia extractiva se traduce automáticamente en prosperidad.
	
	A pesar de la extensa producción académica sobre los efectos económicos de la dependencia extractiva, la evidencia empírica continúa siendo ambigua y, en muchos casos, contradictoria. Algunos estudios encuentran asociaciones negativas robustas entre dependencia de recursos y desempeño económico, mientras que otros destacan efectos neutros o incluso positivos condicionados a la calidad institucional, la acumulación de capital humano o la gestión macroeconómica. Esta heterogeneidad de resultados sugiere que la relación entre recursos naturales y transición estructural no responde a un mecanismo único, sino a configuraciones productivo-institucionales específicas cuya interacción aún no se encuentra plenamente sistematizada en un marco comparativo de largo plazo.
	
	En este contexto, la investigación se propone responder a la siguiente pregunta: ¿bajo qué configuraciones productivo-institucionales la dependencia de recursos naturales no renovables se traduce en trayectorias de diversificación sostenida y acumulación de capacidades en economías intensivas en recursos? El análisis se concentra en el período 1980–2022. Para abordar esta cuestión, el trabajo adopta como eje analítico la noción de transición estructural, entendida como un proceso sostenido de reasignación de recursos hacia actividades de mayor complejidad productiva y tecnológica, acompañado por aprendizaje, innovación y fortalecimiento institucional. Desde esta perspectiva, la dependencia extractiva no constituye un destino predeterminado, sino una condición inicial cuyos efectos dependen de la interacción entre mecanismos productivos (como la enfermedad holandesa y los encadenamientos sectoriales), arreglos institucionales y restricciones macroeconómicas intertemporales.
	
	Empíricamente, la investigación se plantea como un estudio cuantitativo comparativo basado en un panel país–año para economías intensivas en recursos durante el período 1980–2022. La propuesta consiste en estimar modelos econométricos que vinculen indicadores de dependencia extractiva con medidas de transición estructural, incorporando el papel de la calidad institucional y un término de interacción entre recursos e instituciones, junto con variables asociadas a volatilidad externa, dinámica del tipo de cambio real, capital humano, innovación, nivel de desarrollo y restricciones fiscales. La selección final de indicadores y especificaciones dependerá de la disponibilidad y consistencia temporal de los datos, así como de la discusión metodológica con la dirección del trabajo. Este enfoque busca evaluar efectos condicionales —en particular, si la calidad institucional atenúa o refuerza los impactos de la dependencia extractiva— controlando por heterogeneidad no observada entre países y efectos comunes en el tiempo.
	
	La contribución esperada del trabajo radica en ofrecer un marco analítico integrado que permita comprender la dependencia extractiva no como una condición estática, sino como parte de configuraciones productivo-institucionales dinámicas que pueden derivar en trayectorias divergentes de transformación estructural. Al articular mecanismos macroeconómicos, institucionales y productivos dentro de un diseño comparativo de largo plazo, la investigación busca aportar evidencia sistemática sobre las condiciones bajo las cuales las rentas de recursos naturales pueden convertirse en capacidades sostenibles, en lugar de consolidar patrones de especialización primaria. En este sentido, el estudio aspira a contribuir tanto al debate académico sobre desarrollo y recursos naturales como a la discusión de políticas orientadas a la diversificación productiva en economías dependientes de recursos.
	
\end{document}